% Figure: a MEMS pressure-sensor diaphragm. A thin square silicon membrane
% clamped on all edges over an etched cavity, deflecting under pressure, with
% piezoresistors placed at the clamped edge midpoints where the bending
% strain is largest. Cross-section schematic plus a plan of gauge placement.
\begin{figure}[htbp]
\centering
\begin{tikzpicture}[scale=1.0,>=latex]
  % ---- cross-section (left) ----
  \begin{scope}[shift={(0,0)}]
    % silicon frame (thick handle wafer) with an etched cavity
    \fill[enggray!25] (-2.2,-1.1) -- (2.2,-1.1) -- (2.2,0) -- (1.1,0)
       -- (0.8,-0.7) -- (-0.8,-0.7) -- (-1.1,0) -- (-2.2,0) -- cycle;
    \draw[engblack, thick] (-2.2,-1.1) -- (2.2,-1.1) -- (2.2,0) -- (1.1,0)
       -- (0.8,-0.7) -- (-0.8,-0.7) -- (-1.1,0) -- (-2.2,0) -- cycle;
    % the thin diaphragm across the top of the cavity, deflected downward
    \draw[engblue, very thick] (-1.1,0) .. controls (-0.4,-0.28) and (0.4,-0.28) .. (1.1,0);
    \draw[engblue, thin, dashed] (-1.1,0) -- (1.1,0);
    % pressure arrows pushing down on the diaphragm
    \foreach \x in {-0.7,-0.35,0,0.35,0.7}{
      \draw[englime,->] (\x,0.55) -- (\x,0.08);
    }
    \node[englime, font=\scriptsize, anchor=south] at (0,0.55) {pressure $p$};
    % thickness callout of the diaphragm
    \draw[engamber,<->] (1.3,0) -- (1.3,-0.12) node[midway,right,font=\scriptsize]{$h$};
    \node[font=\scriptsize, anchor=north] at (0,-1.25) {cross-section};
    \node[enggray, font=\scriptsize, anchor=center] at (0,-0.5) {cavity};
  \end{scope}
  % ---- plan view (right): gauge placement ----
  \begin{scope}[shift={(5.3,-0.55)}]
    \draw[engblack, thick] (-1.1,-1.1) rectangle (1.1,1.1);
    \draw[engblue, thick, dashed] (-0.75,-0.75) rectangle (0.75,0.75);
    \node[engblue, font=\scriptsize, anchor=south] at (0,0.78) {diaphragm edge};
    % piezoresistors at the four edge midpoints (max bending strain)
    \fill[engred] (-0.75,0) circle (0.06);
    \fill[engred] (0.75,0) circle (0.06);
    \fill[engred] (0,-0.75) circle (0.06);
    \fill[engred] (0,0.75) circle (0.06);
    \node[engred, font=\scriptsize, anchor=west] at (0.85,0) {gauge};
    \node[font=\scriptsize, anchor=north] at (0,-1.25) {plan: gauges at edge midpoints};
  \end{scope}
\end{tikzpicture}
\caption[MEMS pressure-sensor diaphragm]{A micromachined pressure-sensor
diaphragm: a thin silicon membrane, clamped on all four edges over a cavity
etched into the wafer, deflecting under the pressure to be measured. In
cross-section (left) the diaphragm bows into the cavity under pressure $p$; a
membrane a few micrometres thick spanning a side of a fraction of a
millimetre is a clamped plate, and its centre deflection and edge bending
follow the clamped-plate relations of the plate section. The transduction
reads the bending strain rather than the deflection: piezoresistors (red) are
diffused at the midpoints of the clamped edges (right), where the plate
bending strain is largest and where a Wheatstone bridge of four gauges
converts the strain to a voltage linear in pressure over the small-deflection
range. The design lives inside the small-deflection limit of the plate
theory: once the deflection approaches the diaphragm thickness the response
turns nonlinear as membrane stretching takes over, which sets the upper end
of the sensor's usable pressure range.}
\label{fig:vol03ch09:mems-diaphragm}
\end{figure}
